\item \model{MiniMax-M2}

\paragraph*{فلسفه فعال‌سازی حداقلی (\en{Mini Activations})}
سری مدل‌های \model{MiniMax-M2} \cite{minimax_m2_2026} بر مبنای اصل بیشینه‌سازی کارایی پیش‌آموزش طراحی شده‌اند:
\begin{itemize}
    \item \textbf{تعداد کل پارامترها:} ۲۲۹.۹ میلیارد پارامتر (\en{229.9B}).
    \item \textbf{پارامترهای فعال به ازای توکن:} تنها ۹.۸ میلیارد پارامتر (\en{9.8B active}).
    \item \textbf{پیکربندی پایه:} ۶۲ لایه دکودر ترانسفورمر، بعد پنهان ۳۰۷۲، دایره واژگان ۲۰۰,۰۶۴ توکن.
    \item \textbf{طول کانتکست پیش‌آموزش:} ۱۹۲,۰۰۰ توکن (\en{192K context}).
\end{itemize}

\paragraph*{انتخاب راهبردی توجه کامل در برابر رد فرضیه \en{SWA}}
آزمایش‌های پیش‌آموزش گسترده در مقیاس تریلیون توکن بر روی مدل‌های هیبرید توجه پنجره لغزان (\en{Hybrid SWA}) \cite{minimax_m2_2026} نشان داد که الگوهای متناوب توجه محلی باعث زوال شدید توانایی بازیابی اطلاعات در کانتکست‌های بلند و استدلال چندمرحله‌ای می‌شوند (افت امتیاز بازیابی \en{RULER} در طول $128\text{K}$ از ۹۰ به ۷۲). بر این اساس، مدل \model{MiniMax-M2} فرضیه استفاده از لایه‌های تنک پنجره‌ای را کنار گذاشته و از سازوکار توجه متراکم با پرسش‌های گروهی (\en{GQA}) با ۴۸ سر پرسش و ۸ سر کلید-مقدار در تمام لایه‌ها به همراه تعبیه دورانی \en{RoPE} بهره می‌برد.

\paragraph*{متخصصان ریزدانه با گیتینگ سیگموئید و بایاس یادگیرنده}
لایه‌های \en{FFN} شامل ۲۵۶ متخصص ریزدانه هستند که ۸ متخصص به ازای هر توکن گزینش می‌شوند. برخلاف روش‌های سنتی مبتنی بر \en{Softmax}، امتیازدهی از طریق تابع سیگموئید مستقل محاسبه می‌گردد:
\begin{equation}
    s_i = \operatorname{Sigmoid}(W_r x + b_i)
\end{equation}
این ساختار قید مجموع-صفر (\en{Zero-Sum}) را حذف کرده و امکان انتخاب همزمان چندین متخصص برجسته را بدون تداخل مقیاسی فراهم می‌کند. پارامترهای بایاس $b_i$ به طور مشترک با وزن‌های شبکه بهینه‌سازی شده و به صورت ضمنی توازن بار متخصصان را بدون نیاز به توابع اتلاف کمکی سنگین تضمین می‌نمایند.

\paragraph*{پیش‌بینی چندتوکنی (\en{MTP}) و تکامل با کپی وزن‌ها}
در فاز اولیه پیش‌آموزش، مدل به یک ماژول \en{MTP} ($K=1$) مجهز است که با ضریب اتلاف اولیه $0.3$ شروع شده و در فاز زوال به $0.1$ می‌رسد. در فاز زوال آموزش ادامه‌دار، ماژول \en{MTP} به $K=3$ تعمیم می‌یابد. به جای مقداردهی اولیه تصادفی، پارامترهای لایه‌ها مستقیماً از شبکه اصلی کپی می‌شوند:
\begin{equation}
    \theta_{\text{MTP}, k}^{(0)} \leftarrow \theta_{\text{Main Block}}
\end{equation}
این راهبرد انحراف بازنمایی لایه‌ها را به حداقل رسانده و سرعت همگرایی ماژول پیش‌بینی را تا چندین برابر افزایش می‌دهد. فرایند پیش‌آموزش روی ۲۹.۲ تریلیون توکن (۱۹.۹T در فاز نرخ ثابت و ۹.۳T در فاز زوال) با بسط تدریجی طول توالی از $8\text{K}$ به $32\text{K}$ و در نهایت $192\text{K}$ توکن تکمیل گردید.